\documentclass[a4paper,11pt]{article}

% 1. 中文支持 (必须使用 XeLaTeX 编译)
% quiet: 抑制 fontspec 对 TeX Live 自带 Fandol 字体的无害 "Script CJK" 警告
\PassOptionsToPackage{quiet}{fontspec}
% xcolor 选项需在 tikz/ctex 载入 xcolor 之前传入,否则 option clash
\PassOptionsToPackage{usenames,dvipsnames}{xcolor}
\usepackage[UTF8]{ctex}

% 2. 绘图包 (实现蓝色梯形标题)
\usepackage{tikz}
\usetikzlibrary{shapes, backgrounds}

% 3. 布局与格式包
\usepackage{latexsym}
\usepackage[empty]{fullpage}
\usepackage[explicit]{titlesec}
\usepackage{marvosym}
\usepackage{verbatim}
\usepackage{enumitem}
\usepackage[hidelinks]{hyperref}
\usepackage{fancyhdr}
\usepackage{tabularx}
\usepackage{xcolor}

%-------------------------------------------
% 颜色定义
\definecolor{MyDarkBlue}{RGB}{0, 50, 100}

% 页边距设置 (优化后让页面更饱满)
\pagestyle{fancy}
\fancyhf{}
\fancyfoot{}
\renewcommand{\headrulewidth}{0pt}
\renewcommand{\footrulewidth}{0pt}

\addtolength{\oddsidemargin}{-0.6in}
\addtolength{\evensidemargin}{-0.6in}
\addtolength{\textwidth}{1.2in}
\addtolength{\topmargin}{-.7in}
\addtolength{\textheight}{1.4in}

\urlstyle{same}
\raggedbottom
\raggedright
\setlength{\tabcolsep}{0in}

% --- 自定义 TikZ 标题样式 ---
\titleformat{\section}
  {}
  {}
  {0pt}
  {%
    \noindent
    \begin{tikzpicture}[baseline=(t.base)]
    \vspace{-2pt}
      \node[inner xsep=10pt, inner ysep=3pt, text=white, font=\large\bfseries, align=left, anchor=base west] (t) at (0,0) {#1};
      \begin{scope}[on background layer]
         \fill[MyDarkBlue] (t.north west) -- (t.north east) -- ([xshift=8pt]t.south east) -- (t.south west) -- cycle;
         \draw[MyDarkBlue, line width=2pt] ([xshift=-68.30pt]t.south east) -- (\linewidth, 0 |- t.south east);
      \end{scope}
    \end{tikzpicture}%
  }
  []
\titlespacing*{\section}{0pt}{12pt}{6pt}

%-------------------------------------------
% 自定义命令
\newcommand{\resumeItem}[1]{
  \item\small{{#1 \vspace{-1pt}}}
}

\newcommand{\resumeEducationHeading}[4]{
  \vspace{-2pt}\item
    \begin{tabular*}{0.97\textwidth}[t]{l@{\extracolsep{\fill}}r}
      \textbf{#1} & #2 \\
      \textit{\small #3} & \textit{\small #4} \\
    \end{tabular*}\vspace{-5pt}
}

\newcommand{\resumeResearchHeading}[2]{
  \vspace{-2pt}\item
    \begin{tabular*}{0.97\textwidth}[t]{l@{\extracolsep{\fill}}r}
      \textbf{#1} & \textbf{#2} \\
    \end{tabular*}\vspace{-7pt}
}

\newcommand{\resumeProjectHeading}[2]{
    \item
    \begin{tabular*}{0.97\textwidth}{l@{\extracolsep{\fill}}r}
      \small#1 & #2 \\
    \end{tabular*}\vspace{-7pt}
}

\newcommand{\resumeSubHeadingListStart}{\begin{itemize}[leftmargin=0.15in, label={}]}
\newcommand{\resumeSubHeadingListEnd}{\end{itemize}}
\newcommand{\resumeItemListStart}{\begin{itemize}[leftmargin=*, label={\small\textbullet}]}
\newcommand{\resumeItemListEnd}{\end{itemize}\vspace{-2pt}}

%-------------------------------------------
%%%%%%  简历正文  %%%%%%%%%%%%%%%%%%%%%%%%%%%%
\begin{document}

%---------- 个人信息 ----------
\begin{center}
    \textbf{\Huge {\ziju{0.2} 郜一贤}} \\ \vspace{8pt}
    \small
    \href{mailto:yisrealbonaparte@outlook.com}{\underline{yisrealbonaparte@outlook.com}} $|$
    +86 199-3494-2785 $|$
    杭州电子科技大学 $\cdot$ 传播学 $|$
    求职意向：新媒体相关岗位
\end{center}

%----------- 教育背景 -----------
\section{教育背景}
  \resumeSubHeadingListStart
    \resumeEducationHeading
      {杭州电子科技大学}{2023.09 -- 2027.07}
      {本科 \ \textbar \ 传播学}{杭州，浙江}
      \resumeItemListStart
        \resumeItem{\textbf{GPA: 4.5 / 5.0} \ \ $|\ \ $ \textbf{主要荣誉：}优秀学生干部，多次获校一、二、三等奖学金}
        \resumeItem{\textbf{主修课程：}传播学原理、非线性编辑、数据新闻、新媒体内容策略等}
        \resumeItem{\textbf{英语水平：}CET-4 \ \textbar \ 雅思 6.0，具备扎实的跨文化传播与英文文献检索能力}
      \resumeItemListEnd
  \resumeSubHeadingListEnd

%----------- 个人技能 -----------
\section{个人技能}
 \begin{itemize}[leftmargin=0.15in, label={}]
    \small{\item{
     \textbf{设计与剪辑软件}{: 熟练操作 \textbf{Adobe After Effects}、\textbf{Photoshop}、\textbf{达芬奇}、剪映、秀米，具备视音频全链路敏捷后期能力} \\
     \textbf{AI 创作工具}{: 熟练构建 AIGC 工作流，运用 \textbf{ChatGPT}、DeepSeek 进行文本提炼与 Prompt 工程；通过 \textbf{即梦}、Midjourney 快速进行视觉物料概念生成} \\
     \textbf{办公与音频软件}{: 熟练运用 WPS Office 系列（Word、Excel数据看板、PPT汇报）；熟练操作 Audacity 进行音频精修与降噪} \\
     \textbf{核心专业能力}{: 爆款内容网感、新媒体矩阵多渠道分发、标准化视觉排版、跨团队流程协同管理} \\
     \textbf{个人新媒体项目}{: 独立从 0 到 1 运营个人小红书账号，孵化 \textbf{3 篇万级互动优质笔记}，\textbf{单篇最高曝光超 120 万}}
    }}
 \end{itemize}

%----------- 校园经历 -----------
\section{校园经历}
  \resumeSubHeadingListStart

    \resumeResearchHeading
      {杭州电子科技大学校辩论协会 \ \textbar \ 宣传部部长}{2023.09 -- 2025.09}
      \resumeItemListStart
        \resumeItem{\textbf{全周期矩阵运营：}主导校辩协官方微信公众号内容规划与活动全周期运营，独立承接赛事预告、深度专访及赛后复盘全链路内容，**实现单篇推文平均阅读量达 1500+**，展现极高粉丝粘性与外溢传播力。}
        \resumeItem{\textbf{视觉与流程标准化：}规范并推行“选题—三审—视觉排版—数据复盘”标准化漏斗工作流，统一组织视觉视觉VI，**推动官方平台总粉丝量增长至 800+**。}
      \resumeItemListEnd

    \resumeResearchHeading
      {杭州电子科技大学红的家园网络工作室 \ \textbar \ 影视部副部长}{2023.09 -- 2025.09}
      \resumeItemListStart
        \resumeItem{\textbf{校媒视觉供给：}核心对接校官方微博视音频素材供给；通过建立标准化拍摄排期表，统筹大型活动素材采集与后期协同机制，**使团队视觉素材交付周期缩短 40\%**，有力保障校级核心媒体矩阵内容输出稳定性。}
      \resumeItemListEnd

  \resumeSubHeadingListEnd

%----------- 科研竞赛 -----------
\section{科研竞赛}
  \resumeSubHeadingListStart

    \resumeResearchHeading
      {《前沿科学》期刊（国家级刊物） \ \textbar \ 第一作者}{2025}
      \resumeItemListStart
        \resumeItem{以\textbf{第一作者}身份发表学术论文《人工智能时代经典影视剧衍生创作的法律边界研究》（CN11-5568/N），系统剖析 AIGC 传播链路、梳理生成式 AI 技术在二次创作中的著作权风险并提出合规防御应对策略。}
      \resumeItemListEnd

    \resumeResearchHeading
      {品牌规划赛 \ \textbar \ 团队 IP 设计成员}{2025.03 -- 2025.05}
      \resumeItemListStart
        \resumeItem{\textbf{视觉体系构建：}主理民宿品牌“灵溪山居”IP 形象及衍生文创产品设计，原创山水概念形象“溪溪”及配套年轻化视觉VI；**全流程引入 AIGC 工具流**，利用 ChatGPT 提炼 Prompt 矩阵并结合 Midjourney、即梦快速生成视觉草图。}
        \resumeItem{\textbf{多维渠道分发：}协同打通文创周边、短视频矩阵与沉浸式体验馆三维传播链路，创意方案在全校/全赛区脱颖而出，最终斩获 \textbf{\color{MyDarkBlue}国家级二等奖}。}
      \resumeItemListEnd

    \resumeResearchHeading
      {全国大学生广告艺术大赛 \ \textbar \ 核心剪辑 / 演员}{2025.03 -- 2025.06}
      \resumeItemListStart
        \resumeItem{\textbf{视音频敏捷创作：}核心负责营销创客微短剧《UU 跑腿迷案》的后期剪辑与角色演绎；**探索“AI+视听”融合应用**，运用即梦 AI 生成辅助场景画面，结合 AI 配音把控短剧叙事节奏与品牌价值传递，作品最终斩获 \textbf{\color{MyDarkBlue}国赛三等奖}。}
      \resumeItemListEnd

    \resumeResearchHeading
      {EOVS 工创赛 \ \textbar \ 团队财务数据分析}{2024.09 -- 2024.11}
      \resumeItemListStart
        \resumeItem{\textbf{数据建模与分析：}作为虚拟仿真赛项核心成员，独立承接财务与市场仿真数据分析，运用 Excel 优化并建立季度多维经营报表模型，**高效支撑团队连续 8 个季度的运营策略模拟与精细化迭代**，助力团队获 \textbf{\color{MyDarkBlue}省赛三等奖}。}
      \resumeItemListEnd

  \resumeSubHeadingListEnd

%----------- 实习与项目经历 -----------
\section{实习与项目经历}
  \resumeSubHeadingListStart

    \resumeResearchHeading
      {尚岸（杭州）教育信息技术有限责任公司 \ \textbar \ 设计部实习生}{2025.07 -- 2025.08}
      \resumeItemListStart
        \resumeItem{\textbf{数字化资产处理：}独立负责人教版小学、初中数字化教材交互板块音频资产的清洗与导入，熟练操作 \textbf{Audacity} 规范执行音频文件分割、去噪与编号入库，确保音频与书籍内容精准适配，实现零差错高效交付。}
        \resumeItem{\textbf{视觉素材设计：}运用 \textbf{Photoshop} 独立承接小学英语单词卡片设计，输出符合教学大纲的标准化、模块化辅助教学视觉物料。}
      \resumeItemListEnd

    \resumeResearchHeading
      {浙江大学出版社数字教材项目 \ \textbar \ 数字化编辑}{2025.03 -- 2025.07}
      \resumeItemListStart
        \resumeItem{\textbf{内容结构化审校：}独立负责浙大数字教材平台书目迁移与资产录入，精准校对教材元数据，累计完成文史、理工、经管等多学科书目 10 余本的结构化录入，**核对与录入准确率达 95\% 以上**。}
      \resumeItemListEnd

  \resumeSubHeadingListEnd

%----------- 荣誉奖项 -----------
\section{荣誉奖项}
  \resumeSubHeadingListStart
    \resumeProjectHeading
      {\textbf{品牌规划赛} 国家级二等奖（团队 IP 设计核心成员）}{2025}
    \resumeProjectHeading
      {\textbf{全国大学生广告艺术大赛} 国赛三等奖（核心剪辑 / 主演）}{2025}
    \resumeProjectHeading
      {\textbf{EOVS 工创赛} 省赛三等奖（团队财务数据分析）}{2024}
    \resumeProjectHeading
      {\textbf{校奖学金}（多次获得） \ \& \ \textbf{优秀学生干部}}{本科}
    \resumeProjectHeading
      {以 \textbf{第一作者} 身份于国家级期刊《前沿科学》发表学术论文}{2025}
  \resumeSubHeadingListEnd

\end{document}